\documentclass{article}
\usepackage{graphicx}
\usepackage{amsmath}
\usepackage{amssymb}
\usepackage{amsfonts}
\usepackage{hyperref}
\usepackage{exercises}
\usepackage{indentfirst}
\usepackage{circledsteps}
\usepackage{amsthm}
\usepackage{csquotes}
\usepackage{yhmath}
\usepackage{comment}
\usepackage{placeins}
\usepackage{geometry}
\usepackage{hyperref}
\usepackage{tikz}
\newcommand*{\oldneg}{\mathord{\sim}}
\newcommand{\Z}{\mathbb{Z}}
\newcommand{\N}{\mathbb{N}}
\let\oldemptyset\emptyset

\theoremstyle{definition}
\newtheorem{definition}{Definition}
\newtheorem{theorem}{Theorem}[section]
\newtheorem{corollary}{Corollary}[theorem]
\newtheorem{lemma}[theorem]{Lemma}

\setlength{\parindent}{.5in}
\geometry{margin=0.5in}
\setlength{\parskip}{12pt}
\title{Chapter 6 Notes}
\begin{document}
\date{}
\maketitle

\section*{Section 1}

\subsection*{Notes}
Note that $F\restriction A$ does not necessarily equal $A\times F''(A)$. In fact, if 
$F''(A)$ contains more than one element (say it contains at least the two elements $b,c$), then obviously $F\restriction A\ne A\times F''(A)$, as, 
for any element $x\in A$, we would have $\langle x,b\rangle$ and $\langle x,c\rangle$ both in $A\times F''(A)$, 
but there can only be one $y$ for which $\langle x,y\rangle\in F\restriction A$.

It is remarked that:\textquote{If $\phi_1$ and $\phi_2$ are functions from a class $A$ onto a class $B$ and if $\phi_1\phi_2^{-1}$ is an identity 
function, then obviously $\phi_1=\phi_2$.} Here is why this is. Now, note that this assumes that $\phi_2^{-1}$, is, in fact, a function. Then, for 
every element $y$ in the range of $\phi_2$, we have $\phi_2^{-1}(y)=x$ for some $x$ in the domain of $\phi_2$, and 
$\phi_1\phi_2^{-1}(y)=\phi_1(x)=y$. But also, $\phi_2(x)=y$. Thus $\phi_1(x)=\phi_2(x)$ for every $x$ in the domain of $\phi_2$, and 
because they have the same domains, the functions are the same.

Note that if a class $A$ is well ordered under $\subseteq$, then it is also properly well ordered under $\subseteq$. The proof of this is the same as 
the proof of Proposition 1.3.

\subsection*{Exercise 1.1}
Let $\mathcal{N}$ be a nest of functions.

(1): We first show that $\cup\mathcal{N}$ is function. Obviously $\cup\mathcal{N}$ is a relation. Now, suppose we 
had $\langle x,y\rangle$ and $\langle x,z\rangle$ both in $\cup\mathcal{N}$ for some $y\ne z$. Then $x$ is in the domain of some 
function $f$. Without loss of generality, suppose $\langle x,y\rangle\in f$. Then $\langle x,z\rangle\notin f$. Now, this would mean 
that $\langle x,z\rangle$ was in some other function $g$. But, because $\mathcal{N}$ is a nest, either $f\subseteq g$ or $g\subseteq f$. Without 
loss of generality, suppose $f\subseteq g$. Then we would have $g(x)=y$ and $g(x)=z$, a contradiction to $g$ being a function. Hence $\cup\mathcal{N}$
must be a function.

(2): We now show that the domain of $\cup\mathcal{N}N$ is the union of the class of domains of the elements of $\mathcal{N}$ (call this $D$). Suppose 
$x\in\textrm{Dom}(\cup\mathcal{N})$. Then, for some $y$ and some function $f$, $\langle x,y\rangle\in f$, meaning $x\in D$. On the other hand, suppose 
$x\in D$. Then, for some $y$ and some function $f$, $langle x,y\rangle\in f$, meaning $x\in\textrm{Dom}(\cup\mathcal{N})$. Thus $\textrm{Dom}(\cup\mathcal{N})=D$. 
Similar reasoning shows $\textrm{Ran}(\cup\mathcal{N})$ is the union of the class of ranges of the elements of $\mathcal{N}$.

(3): Lastly, suppose that every element of $\mathcal{N}$ was 1-1. Now, suppose $x$ and $y$ are in the domain of $\cup\mathcal{N}$, and that $x\ne y$, and 
for some $v$ and $w$, $\langle x,v\rangle\in\cup\mathcal{N}$ and $\langle y,w\rangle\in\cup\mathcal{N}$. 
Then, for some functions $f$ and $g$ in $\mathcal{N}$, $\langle x,v\rangle\in f$ and $\langle y,w\rangle\in g$. 
Without loss of generality, suppose that $f\subseteq g$. Then $g$ is 1-1, meaning $v\ne w$. Hence $\cup\mathcal{N}$ is 1-1.

\subsection*{Exercise 1.2}
First, let us note that $On$ is not a set. For if it were, it would then be a transitive set of ordinals, and thus would be an ordinal. Then that would mean 
$On\in On$, contrary to Theorem 1.8 (of Chapter 5). Thus $On$ is not a set.

Now, let $A$ be the class $On\cup\{\{\{\emptyset\}\}\}$. Thus $x\in A$ iff either $x$ is an ordinal, or $x=\{\{\emptyset\}\}$. Note that 
$\{\{\emptyset\}\}$ is not an ordinal, as if it were, then it would be transitive, which would imply that $\emptyset\in\{\{\emptyset\}\}$, which 
is not the case. Hence $On$ is a proper subclass of $A$.

Then define the relation $x\le y$ as follows: $x\le y$ iff $x$ and $y$ are both ordinals and $x\subseteq y$, or $y=\{\{\emptyset\}\}$. Note that this defintion 
implies that $\{\{\emptyset\}\}\le y$ only holds if $y=\{\{\emptyset\}\}$. We first show that 
$\le$ is a well ordering of $A$.

$O_1$: If $x$ is an ordinal, then $x\subseteq x$, so that $x\le x$ holds. If $x=\{\{\emptyset\}\}$, then $x\le x$ by definition.

$O_2$: Suppose $(x\le y)\wedge (y\le z)$. If all of $x,y$ and $z$ are ordinals, then obviously $x\le z$. In fact, the only one of $x,y$ and $z$ which 
could be $\{\{\emptyset\}\}$ is $z$. In that case, $x\le z$ holds by definition.

$O_3$: Suppose $x\le y$ and $y\le x$. Then either $x=y=\{\{\emptyset\}\}$, or both $x$ and $y$ are ordinals, in which case 
$x\subseteq y$ and $y\subseteq x$, so that $x=y$. Thus $x=y$.

$O_4$: If $x$ and $y$ are both ordinals, then of course either $x\le y$ or $y\le x$. If $y=\{\{\emptyset\}\}$, then $x\le y$ (similarly with $x$).

$O_5$: For any subclass $B$ of $A$, either $\{\{\emptyset\}\}\in B$ or not. If it is, then, $\{\{\emptyset\}\}$ is the least element of $B$ 
(as $x\le\{\{\emptyset\}\}$ holds for all elements in $A$). If it is not in $B$, then $B$ is a subclass of $On$, hence (because $On$ is well 
ordered under inclusion), $B$ has a least element. 

Thus $\le$ is a well ordering of $A$. Now, $On$ is a proper lower section of $A$: It is a proper subclass of $A$, and 
$x<\{\{\emptyset\}\}$ for all elements in $On$. But $On$ is not a set. Thus we have an example of a well ordering which 
is not a proper well ordering.

\end{document}
